\documentclass[11pt,a4paper]{article}

% -----------------------------------------------------------------------
% Packages
% -----------------------------------------------------------------------
\usepackage{amsmath,amsthm,amssymb,mathtools}
\usepackage[T1]{fontenc}
\usepackage{lmodern}
\usepackage{microtype}
\usepackage{geometry}
\usepackage{booktabs}
\usepackage{array}
\usepackage{enumitem}
\usepackage{hyperref}
\usepackage{url}
\usepackage{xcolor}

\geometry{margin=1in}

\hypersetup{
  colorlinks=true,
  linkcolor=blue!70!black,
  citecolor=blue!70!black,
  urlcolor=blue!70!black,
}

% -----------------------------------------------------------------------
% Theorem environments
% -----------------------------------------------------------------------
\theoremstyle{plain}
\newtheorem{theorem}{Theorem}[section]
\newtheorem{lemma}[theorem]{Lemma}
\newtheorem{proposition}[theorem]{Proposition}
\newtheorem{corollary}[theorem]{Corollary}

\theoremstyle{definition}
\newtheorem{definition}[theorem]{Definition}

\theoremstyle{remark}
\newtheorem{remark}[theorem]{Remark}
\newtheorem{observation}[theorem]{Observation}
\newtheorem{conjecture}[theorem]{Conjecture}

% -----------------------------------------------------------------------
% Macros
% -----------------------------------------------------------------------
\newcommand{\N}{\mathbb{N}}
\newcommand{\Z}{\mathbb{Z}}
\newcommand{\floor}[1]{\left\lfloor #1 \right\rfloor}
\newcommand{\ceil}[1]{\left\lceil #1 \right\rceil}
\newcommand{\val}{\operatorname{val}}
\newcommand{\ones}{N_1}
\newcommand{\zeros}{N_0}
\newcommand{\Da}{\Delta a}
\newcommand{\DDa}{\Delta^2 a}

% -----------------------------------------------------------------------
% Title
% -----------------------------------------------------------------------
\title{On the maximum number of distinct integers\\
       representable by substrings of a binary string}

\author{Ben Clare}
\date{\today}

% =======================================================================
\begin{document}
% =======================================================================

\maketitle

% -----------------------------------------------------------------------
\begin{abstract}
% -----------------------------------------------------------------------
For a binary string $s$ of length~$n$, let $f(s)$ denote the number of
distinct non-negative integers obtained by reading contiguous substrings
of~$s$ as binary numerals, with leading zeros ignored, and set
$a(n) = \max_{|s|=n} f(s)$.  This is OEIS sequence A112509.

We prove that
$\lim_{n\to\infty} a(n)/n^2 = \frac{1}{2}$; more precisely, for all
$n \ge 36$,
\[
  \tfrac{1}{2}\,n^2 - \tfrac{13}{4}\,n^{3/2}
  \;\le\; a(n)
  \;\le\; \tfrac{n(n+1)}{2},
\]
so $a(n)/n^2 = \frac{1}{2} + O(n^{-1/2})$.
We also construct explicit near-optimal strings of ones-density
$1-O(n^{-1/2})$ and prove that every optimal string has ones-density
$1-O(n^{-1/4})$.
The sequence is convex:
$a(n{+}1) - 2a(n) + a(n{-}1) \ge 0$ for all $n \ge 2$, and its second
differences satisfy $\Delta^2 a(n)=O(n^{3/4})$.
Finally, we prove that any two consecutive zero-separators of equal
length force
$\Theta(\min(b_i,b_{i+1})\cdot \min(b_{i+1},b_{i+2}))$ duplicate
substring values, giving a quantitative obstruction to repeated
adjacent separator lengths.

A branch-and-bound algorithm founded on these bounds extends the
certified table from Fuller's A$^*$ search for $n \le 80$ to
$n \le 111$.  An explicitly scored string at $n = 2 \times 10^9$ gives
the certified lower bound $a(n)/n^2 \ge 0.49996$.
\end{abstract}

\tableofcontents

% =======================================================================
\section{Introduction}\label{sec:intro}
% =======================================================================

Given a binary string, every contiguous substring can be read as the
binary representation of a non-negative integer (with leading zeros
suppressed: $\texttt{01} = \texttt{1} = 1$, and $\texttt{0} = 0$).
The question of maximising the number of distinct values simultaneously
representable in a single $n$-bit string was motivated by a competition
problem and appears as OEIS A112509~\cite{oeis}.

\begin{definition}\label{def:main}
  For $s \in \{0,1\}^n$, let $V(s)$ denote the set of distinct
  non-negative integers representable by contiguous substrings of~$s$,
  and set $f(s) = |V(s)|$.  Define
  \[
    a(n) = \max_{s \in \{0,1\}^n} f(s).
  \]
  We use the convention $a(0) = 0$.
\end{definition}

\begin{lemma}[Leading-bit reduction]\label{lem:leading-bit}
  For every $n \ge 2$, every optimal string for $a(n)$ begins
  with~$\mathtt{1}$.
\end{lemma}

\begin{proof}
  Let $s = \mathtt{0}u$ have length~$n$, and set
  $t = \mathtt{1}u$.  Every substring of~$s$ not beginning at the
  first position is also a substring of~$t$.  A substring of~$s$
  beginning at the first position has the form $\mathtt{0}v$, where
  $v$ is a prefix of~$u$; if $v$ is nonempty, then
  $\val(\mathtt{0}v) = \val(v)$ and $v$ occurs in~$t$ beginning at
  the second position.  Thus changing the first bit from
  $\mathtt{0}$ to $\mathtt{1}$ cannot lose any positive value.

  If $u$ contains a zero, the value~$0$ is still represented in~$t$,
  and the full word $t$ represents a positive integer whose standard
  representation has length~$n$; this value is not represented in~$s$.
  Hence $f(t) > f(s)$.

  It remains only to consider the case in which $u$ contains no zero,
  i.e. $u = \mathtt{1}^{n-1}$.
  If $u = \mathtt{1}^{n-1}$, then
  $V(s) = \{0,1,3,\ldots,2^{n-1}-1\}$ has size~$n$.  But
  the string $\mathtt{1}^{x}\mathtt{0}^{y}$ with
  $x=\floor{n/2}$ and $y=\ceil{n/2}$ has at least $x$ all-one
  values, one zero value, and $xy$ mixed values
  $\val(\mathtt{1}^a\mathtt{0}^b)$ distinguished by their odd parts
  and $2$-adic valuations.  Thus it has
  $xy+x+1>n$ distinct values for every $n\ge2$.  Hence
  such an~$s$ is not optimal either.
\end{proof}

The first twenty values of the sequence are
\[
  1,\; 3,\; 5,\; 7,\; 10,\; 13,\; 17,\; 22,\; 27,\; 33,\;
  40,\; 47,\; 55,\; 64,\; 73,\; 83,\; 94,\; 106,\; 118,\; 131.
\]

Prior to the present work, values for $n \le 80$ were certified by
M.~Fuller using an A$^*$ search and updated in the OEIS log in 2025.
No asymptotic results were available in the literature.

\subsection*{Main theorem}

The following theorem collects the main rigorous contributions of the
paper.  Heuristic observations not covered by this statement are
separated later in Section~\ref{sec:observations}.

\begin{theorem}[Main results]\label{thm:main-results}
  The sequence $a(n)$ has the following properties.
  \begin{enumerate}[label=(\roman*),nosep]
    \item The asymptotic limit exists and equals
      \[
        \lim_{n\to\infty} \frac{a(n)}{n^2} = \frac12.
      \]
      More precisely, for every $n \ge 36$,
      \[
        \frac{n^2}{2} - \frac{13}{4}n^{3/2}
        \le a(n) \le \frac{n(n+1)}{2}.
      \]
    \item There are strings $s_n$ with
      $f(s_n) \ge n^2/2 - O(n^{3/2})$ and
      proportion of ones $1 - O(n^{-1/2})$.  Conversely, every
      optimal string contains $O(n^{3/4})$ zeros.
    \item The sequence is convex:
      $a(n+1)-2a(n)+a(n-1) \ge 0$ for every $n \ge 2$.
      Moreover $\DDa(n)=O(n^{3/4})$.
    \item The values of $a(n)$, and the complete set of optimal
      strings, are certified for $1 \le n \le 111$: Fuller's A$^*$ search
      covers $n \le 80$, and the branch-and-bound computation in this
      work covers $81 \le n \le 111$.
    \item If two consecutive zero-separators of a binary string have
      equal length~$s$, and the three adjacent one-block lengths are
      $b_i,b_{i+1},b_{i+2}$, then they force
      $\min(b_i,b_{i+1})\min(b_{i+1},b_{i+2})$ duplicate substring
      values of the form $\val(1^j0^s1^k)$.
  \end{enumerate}
\end{theorem}

\begin{proof}
  Part~(i) is Theorem~\ref{thm:limit}, with the explicit bounds from
  Theorems~\ref{thm:upper} and~\ref{thm:quant-sharp}.  Part~(ii) is
  Proposition~\ref{prop:near-opt-density} and
  Proposition~\ref{prop:optimal-density}.  Part~(iii) is
  Theorem~\ref{thm:convex} together with
  Corollary~\ref{cor:dda-sublinear}.  Part~(iv) is
  Theorem~\ref{thm:certified-computation}.  Part~(v) is
  Theorem~\ref{thm:sep2-collisions}.
\end{proof}

\subsection*{Proof strategy and contributions}

\begin{enumerate}[label=(\roman*)]
  \item The upper bound is the elementary count of possible
    leading-$\mathtt{1}$ substrings; the matching lower bound is an
    explicit family of strings whose internal one-block lengths form a
    strictly decreasing partition.  This avoids probabilistic existence
    arguments and gives explicit constants.
  \item The structural bounds convert the deficit
    $n(n+1)/2-a(n)$ into restrictions on zeros and run lengths.  This
    gives density bounds for optimal strings and explains why long
    monochromatic runs are forced to be sublinear.
  \item Convexity is proved by comparing an optimal string with its
    prefix and suffix.  A zero-dichotomy lemma isolates the only
    obstruction caused by the value~$0$.
  \item The exact computation is certified by combining Fuller's A$^*$
    search for $n \le 80$ with a branch-and-bound search for
    $81 \le n \le 111$.  The branch-and-bound pruning rule is a
    theorem: the suffix can add at most a certified residual
    contribution plus a controlled number of crossing values.  The
    suffix-automaton refinement only strengthens this upper bound.
  \item The separator theorem is stated as a proved local obstruction.
    Stronger separator patterns observed in the data are recorded
    separately as observations and conjectures.
\end{enumerate}

The accompanying data also record the related OEIS companion sequences
obtained from optimal strings, including the smallest and largest
optimal strings and the number of optima.

% =======================================================================
\section{Definitions and notation}\label{sec:definitions}
% =======================================================================

Throughout, binary strings are finite sequences over $\{0,1\}$.  The
\emph{binary value} of $w = w_1 \cdots w_k \in \{0,1\}^k$ is
\[
  \val(w) = \sum_{i=1}^{k} w_i \, 2^{k-i}.
\]
A \emph{substring} of $s = s_1 \cdots s_n$ is any $s_i s_{i+1} \cdots
s_j$ with $1 \le i \le j \le n$; we write $s[i..j]$ for this
substring.  The value set is
$V(s) = \{\val(s[i..j]) : 1 \le i \le j \le n\}$.

The \emph{run-length encoding} (RLE) of $s$ is its decomposition into
maximal constant blocks.  When $s$ begins with~$\mathtt{1}$, we write
$s = 1^{b_1} 0^{c_1} 1^{b_2} 0^{c_2} \cdots$, where the $b_i$ are
one-block lengths and the $c_i$ are zero-block lengths.  The zero-block
lengths $c_i$, including a possible terminal zero-block, are called
\emph{separators}.  Write $\ones(s)$ and $\zeros(s)$ for the total
number of ones and zeros in~$s$.

If $s$ has $1$-runs of lengths $\ell_1,\ldots,\ell_p$ and $0$-runs of
lengths $m_1,\ldots,m_q$, set
$L(s) = \max_t \ell_t$ and $K(s) = \max_t m_t$, with the convention
that the maximum of an empty set is~$0$.

The \emph{first difference} is $\Da(n) = a(n) - a(n{-}1)$ and, for
$n \ge 2$, the \emph{second difference} is
$\DDa(n) = \Da(n{+}1)-\Da(n) = a(n{+}1) - 2a(n) + a(n{-}1)$.

% =======================================================================
\section{Upper bounds}\label{sec:upper}
% =======================================================================

\begin{theorem}[Trivial upper bound]\label{thm:upper}
  For every $n \ge 1$, $\;a(n) \le n(n{+}1)/2$.
  In particular, $\limsup_{n\to\infty} a(n)/n^2 \le \frac{1}{2}$.
\end{theorem}

\begin{proof}
  Every distinct positive integer in $V(s)$ has a unique standard
  binary representation beginning with~$\mathtt{1}$; distinct
  positive values therefore correspond to distinct
  leading-$\mathtt{1}$ substrings of~$s$.

  If $s$ contains at least one zero, then $0 \in V(s)$, represented
  only by substrings beginning with~$\mathtt{0}$.  At least one of
  the $n(n{+}1)/2$ substrings begins with~$\mathtt{0}$, so the
  number of leading-$\mathtt{1}$ substrings is at most
  $n(n{+}1)/2 - 1$, giving $f(s) \le n(n{+}1)/2$.

  If $s = \mathtt{1}^n$, then $0 \notin V(s)$ and
  $f(s) = n \le n(n{+}1)/2$.
\end{proof}

\begin{remark}\label{rem:tight-upper}
  The bound is tight at $n = 2$:
  $V(\mathtt{10}) = \{0, 1, 2\}$, so $a(2) = 3 = 2 \cdot 3/2$.
\end{remark}

\begin{theorem}[Length-stratified upper bound]\label{thm:length-strat}
  Let $K(n) = \max\{\ell \ge 1 : 2^{\ell-1} \le n - \ell + 1\}$.  Then
  \[
    a(n) \;\le\; 2^{K(n)} + \binom{n - K(n) + 1}{2}.
  \]
\end{theorem}

\begin{proof}
  Stratify the distinct positive values by the length~$\ell$ of their
  standard binary representation.  The number $D_\ell$ of distinct
  positive values representable by length-$\ell$ substrings of~$s$
  satisfies $D_\ell \le \min(2^{\ell-1},\, n - \ell + 1)$: the first
  bound counts binary strings of length~$\ell$ starting
  with~$\mathtt{1}$; the second counts length-$\ell$ windows in~$s$.
  Adding at most~$1$ for the value~$0$:
  \[
    a(n) \le 1 + \sum_{\ell=1}^{n} \min(2^{\ell-1},\, n - \ell + 1).
  \]
  By definition of~$K = K(n)$, the minimum equals $2^{\ell-1}$ for
  $\ell \le K$ and $n - \ell + 1$ for $\ell > K$.  Hence
  \[
    1 + \sum_{\ell=1}^{K}(2^{\ell-1})
      + \sum_{\ell=K+1}^{n}(n - \ell + 1)
    = 2^{K} + \binom{n - K + 1}{2}. \qedhere
  \]
\end{proof}

\begin{remark}
  Since $K(n) = \log_2 n + O(\log\log n)$, the bound equals
  $n^2/2 - O(n \log n)$, which is strictly sharper than
  Theorem~\ref{thm:upper} for every $n \ge 3$.
\end{remark}

% =======================================================================
\section{The elementary lower bound}\label{sec:lower}
% =======================================================================

\begin{theorem}[Elementary lower bound]\label{thm:lower}
  For all $n \ge 2$,
  \[
    a(n) \;\ge\; \floor{n/2}\,\ceil{n/2} + \floor{n/2} + 1
         \;\ge\; \frac{n^2}{4} + \frac{n}{2}.
  \]
  In particular, $\liminf_{n\to\infty} a(n)/n^2 \ge \frac{1}{4}$.
\end{theorem}

\begin{proof}
  Set $x = \floor{n/2}$, $y = \ceil{n/2}$, and consider
  $s = \mathtt{1}^x \mathtt{0}^y$.  Its substrings fall into three
  disjoint families:
  \begin{enumerate}[label=(\alph*)]
    \item \textbf{All-ones:} $s[i..j]$ with $j \le x$, giving values
      $2^k - 1$ for $k = 1,\ldots,x$: these are $x$ distinct odd
      positive integers.
    \item \textbf{All-zeros:} $s[i..j]$ with $i > x$, all having
      value~$0$: one distinct value.
    \item \textbf{Mixed:} $s[i..j]$ with $i \le x < j$, of the form
      $\mathtt{1}^a \mathtt{0}^b$ with value $(2^a - 1) \cdot 2^b$,
      where $a \in \{1,\ldots,x\}$ and $b \in \{1,\ldots,y\}$.
      These $xy$ values are pairwise distinct: for fixed~$b$,
      different~$a$ give different odd parts $(2^a - 1)$; for
      different~$b$, the $2$-adic valuations differ.  Moreover, all
      mixed values are even (disjoint from all-ones) and positive
      (disjoint from zero).
  \end{enumerate}
  Thus $f(s) = x + 1 + xy = x(y + 1) + 1$.  With $x = \floor{n/2}$
  and $y = \ceil{n/2}$: for even $n = 2m$,
  $f(s) = m(m+1) + 1 = n^2/4 + n/2 + 1$; for odd $n = 2m+1$,
  $f(s) = m(m+2) + 1 = (n+1)^2/4 > n^2/4 + n/2$.
\end{proof}

% =======================================================================
\section{The asymptotic limit}\label{sec:limit}
% =======================================================================

The elementary lower bound $a(n) \ge n^2/4 + O(n)$ and the upper bound
$a(n) \le n(n{+}1)/2$ leave the leading constant of $a(n)/n^2$
undetermined within $[1/4, 1/2]$.  We now show it equals~$1/2$ by
constructing binary strings with $f(s) \ge n^2/2 - O(n^{3/2})$.

\subsection{Strictly decreasing partitions}

\begin{lemma}[Decreasing partition]\label{lem:partition}
  Let $k \ge 2$ and $M \ge k(k{+}1)/2$ (equivalently,
  $M/k \ge (k{+}1)/2$).  There exist strictly
  decreasing positive integers $b_1 > b_2 > \cdots > b_k \ge 1$
  with $\sum_{i=1}^k b_i = M$.  Moreover, they may be chosen with
  $|b_i - M/k| \le (k{+}1)/2$ for all~$i$.
\end{lemma}

\begin{proof}
  Set $q = \lfloor M/k - (k{-}1)/2 \rfloor$ and
  $r = M - kq - k(k{-}1)/2$.

  \textit{Step~1: $0 \le r < k$.}
  From the floor definition, $q \le M/k - (k{-}1)/2 < q + 1$.
  Multiplying by~$k$: $kq \le M - k(k{-}1)/2 < k(q{+}1)$,
  so $0 \le r < k$.

  \textit{Step~2: Construction.}  Define
  \[
    b_i =
    \begin{cases}
      q + k - i + 1, & 1 \le i \le r, \\
      q + k - i,     & r < i \le k.
    \end{cases}
  \]

  \textit{Step~3: Strict decrease.}
  For consecutive $i, i{+}1$ both in the first case:
  $b_i - b_{i+1} = 1 > 0$.
  At the transition $i = r$ (if $r \ge 1$):
  $b_r - b_{r+1} = (q + k - r + 1) - (q + k - r - 1) = 2 > 0$.
  Both in the second case: $b_i - b_{i+1} = 1 > 0$.

  \textit{Step~4: Sum equals~$M$.}
  \begin{align*}
    \sum_{i=1}^k b_i
    &= \sum_{i=1}^{r}(q + k - i + 1) + \sum_{i=r+1}^{k}(q + k - i) \\
    &= \sum_{i=1}^{k}(q + k - i) + r
    = kq + \frac{k(k{-}1)}{2} + r = M.
  \end{align*}

  \textit{Step~5: Positivity.}
  From $M \ge k(k{+}1)/2$, we get $M/k - (k{-}1)/2 \ge 1$, so
  $q \ge 1$ and $b_k = q \ge 1$.

  \textit{Step~6: Proximity.}
  All $b_i \in [q,\, q + k]$ and
  $M/k \in [q + (k{-}1)/2,\; q + (k{+}1)/2)$, so
  $|b_i - M/k| \le (k{+}1)/2$.
\end{proof}

\subsection{The explicit construction}

Fix $n \ge 64$ and set
\begin{equation}\label{eq:params}
  k = \floor{\sqrt{n}/4}, \qquad M = n - (k - 1).
\end{equation}
Since $n \ge 64$, we have $k \ge 2$.  Moreover,
$k - 1 < \sqrt{n}/4 \le n/2$, so $M > n/2$; and since
$k(k{+}1) \le 2k^2 \le n/8 \le n$, we have
$M \ge n/2 \ge k(k{+}1)/2$, so Lemma~\ref{lem:partition} applies.

Choose $b_1 > b_2 > \cdots > b_k \ge 1$ with $\sum b_i = M$ and
$|b_i - M/k| \le k$, as guaranteed by Lemma~\ref{lem:partition}.
Define the binary string
\begin{equation}\label{eq:string}
  s = \underbrace{1^{b_1}}_{\text{block~1}}\,0\,
      \underbrace{1^{b_2}}_{\text{block~2}}\,0\,
      \cdots\,0\,
      \underbrace{1^{b_k}}_{\text{block~$k$}},
\end{equation}
consisting of $k$ one-blocks separated by $k - 1$ single zeros.  Its
length is $\sum b_i + (k{-}1) = M + (k{-}1) = n$.

\subsection{An injective family of substrings}

For integers $i, j, \alpha, \beta$ satisfying
\begin{equation}\label{eq:range}
  1 \le i, \quad j \ge i + 2, \quad j \le k, \quad
  1 \le \alpha \le b_i, \quad 1 \le \beta \le b_j,
\end{equation}
define $u_{i,j,\alpha,\beta}$ to be the substring of~$s$ beginning
$\alpha$ positions before the end of block~$i$ and ending after the
first $\beta$ ones of block~$j$.  Explicitly, its run-length encoding is
\begin{equation}\label{eq:u}
  \begin{aligned}
  u_{i,j,\alpha,\beta}
  &= 1^{\alpha}\,0\,1^{b_{i+1}}\,0\,\cdots\,0\,
     1^{b_{j-1}}\,0\,1^{\beta}.
  \end{aligned}
\end{equation}
The condition $j \ge i + 2$ ensures at least one full interior block
$b_{i+1},\ldots,b_{j-1}$.  Every $u_{i,j,\alpha,\beta}$ starts and
ends with~$\mathtt{1}$, so it is the standard binary representation of
a positive integer.

\begin{lemma}[Injectivity]\label{lem:inject}
  The map $(i,j,\alpha,\beta) \mapsto u_{i,j,\alpha,\beta}$ is
  injective on the range~\eqref{eq:range}.
\end{lemma}

\begin{proof}
  Suppose $u_{i,j,\alpha,\beta} = u_{i',j',\alpha',\beta'}$ as
  binary words.  The run-length encoding of~\eqref{eq:u} is
  \[
    (\alpha, 1, b_{i+1}, 1, \ldots, b_{j-1}, 1, \beta).
  \]
  Matching entries gives $\alpha = \alpha'$, $\beta = \beta'$, and
  $(b_{i+1},\ldots,b_{j-1}) = (b_{i'+1},\ldots,b_{j'-1})$.  Since
  $b_1 > \cdots > b_k$ is strictly decreasing, every consecutive
  subsequence occurs in exactly one position, so $i = i'$ and
  $j = j'$.
\end{proof}

\subsection{Quantitative lower bound}

\begin{theorem}[Quantitative lower bound]\label{thm:quant}
  For every $n \ge 64$,
  $\;a(n) \ge n^2/2 - 15\,n^{3/2}$.
\end{theorem}

\begin{proof}
  By Lemma~\ref{lem:inject}, the family
  $\{u_{i,j,\alpha,\beta}\}$ consists of pairwise distinct positive
  integers (each starts and ends with~$\mathtt{1}$), so
  \begin{equation}\label{eq:f-lower}
    a(n) \ge f(s) \ge \sum_{\substack{1 \le i < j \le k \\ j \ge i+2}} b_i b_j.
  \end{equation}
  Expanding:
  \begin{equation}\label{eq:expand}
    \sum_{\substack{i < j \\ j \ge i+2}} b_i b_j
    = \frac{(\sum b_i)^2 - \sum b_i^2}{2}
      - \sum_{i=1}^{k-1} b_i b_{i+1}.
  \end{equation}

  \textit{Step~1.}  Since $\sum b_i = M = n - k + 1$ and
  $k \le \sqrt{n}/4$:
  \[
    \frac{M^2}{2} \ge \frac{(n - \sqrt{n}/4)^2}{2}
    \ge \frac{n^2}{2} - \frac{n^{3/2}}{4}.
  \]

  \textit{Step~2.}  Writing $b_i = M/k + d_i$ with $|d_i| \le k$:
  $\sum b_i^2 = M^2/k + \sum d_i^2 \le M^2/k + k^3$.
  Since $k \ge \sqrt{n}/8$ (as $\lfloor x \rfloor \ge x/2$ for
  $x \ge 2$), $M^2/k \le n^2/(\sqrt{n}/8) = 8n^{3/2}$ and
  $k^3 \le n^{3/2}/64$, giving
  \begin{equation}\label{eq:sumsq}
    \sum_{i=1}^k b_i^2 < 9n^{3/2}.
  \end{equation}

  \textit{Step~3.}  By the AM--GM inequality,
  $b_i b_{i+1} \le (b_i^2 + b_{i+1}^2)/2$, so
  \begin{equation}\label{eq:adj}
    \sum_{i=1}^{k-1} b_i b_{i+1} \le \sum_{i=1}^k b_i^2 < 9n^{3/2}.
  \end{equation}

  \textit{Conclusion.}  Substituting into~\eqref{eq:expand}:
  \[
    a(n) \ge \frac{n^2}{2} - \frac{n^{3/2}}{4}
    - \frac{9n^{3/2}}{2} - 9n^{3/2}
    = \frac{n^2}{2} - \frac{55}{4}\,n^{3/2}
    > \frac{n^2}{2} - 15n^{3/2}. \qedhere
  \]
\end{proof}

\begin{theorem}[Limit]\label{thm:limit}
  $\displaystyle\lim_{n \to \infty} \frac{a(n)}{n^2} = \frac{1}{2}$.
\end{theorem}

\begin{proof}
  Theorem~\ref{thm:upper} gives
  $\limsup a(n)/n^2 \le 1/2$.  Theorem~\ref{thm:quant} gives
  $a(n)/n^2 \ge 1/2 - 15/\sqrt{n} \to 1/2$.  The squeeze theorem
  gives the result.
\end{proof}

\subsection{A sharper quantitative bound}

The constant $15$ in Theorem~\ref{thm:quant} is an artifact of the
conservative parameter choice $k = \lfloor\sqrt{n}/4\rfloor$.  By
choosing $k = \lceil\sqrt{n}\,\rceil$ instead, we reduce the constant
by a factor exceeding~$4$.

\begin{theorem}[Sharper quantitative lower bound]\label{thm:quant-sharp}
  For every $n \ge 36$,
  $\;a(n) \ge n^2/2 - (13/4)\,n^{3/2}$.
\end{theorem}

\begin{proof}
  The construction~\eqref{eq:string} and
  Lemmas~\ref{lem:partition}--\ref{lem:inject} are unchanged; only
  the parameter choice differs.  Set $k = \ceil{\sqrt{n}}$ and
  $M = n - k + 1$.  Then $\sqrt{n} \le k \le \sqrt{n} + 1$.

  \textit{Lemma~\ref{lem:partition} applies for $n \ge 36$.}
  We need $M \ge k(k{+}1)/2$.  Since $k \le \sqrt{n} + 1$, this
  requires $n - \sqrt{n} \ge \frac{1}{2}(\sqrt{n}+1)(\sqrt{n}+2)
  = n/2 + 3\sqrt{n}/2 + 1$, i.e., $n/2 \ge 5\sqrt{n}/2 + 1$, which holds for
  $\sqrt{n} \ge 6$, i.e., $n \ge 36$.

  \textit{Variance bound.}  By Lemma~\ref{lem:partition}, each $b_i$
  lies in $[q, q{+}k]$, giving
  \begin{equation}\label{eq:sumsq-sharp}
    \sum_{i=1}^k b_i^2
    \le \frac{M^2}{k} + \frac{k(k{+}1)^2}{4}.
  \end{equation}
  By AM--GM,
  $\sum_{i=1}^{k-1} b_i b_{i+1} \le \sum b_i^2$.
  Substituting into~\eqref{eq:expand}:
  \[
    a(n) \ge \frac{M^2}{2}
      - \frac{3M^2}{2k} - \frac{3k(k{+}1)^2}{8}.
  \]

  \textit{Term-by-term estimates} using
  $\sqrt{n} \le k \le \sqrt{n} + 1$ and $M = n - k + 1$:
  \begin{enumerate}[label=(\roman*),nosep]
    \item $\dfrac{M^2}{2} \ge \dfrac{(n - \sqrt{n})^2}{2}
      = \dfrac{n^2}{2} - n^{3/2} + \dfrac{n}{2}$.
    \item $\dfrac{3M^2}{2k} \le \dfrac{3n^2}{2\sqrt{n}} = \dfrac{3}{2}\,n^{3/2}$.
    \item Expanding $(\sqrt{n}+1)(\sqrt{n}+2)^2 = n^{3/2} + 5n + 8\sqrt{n} + 4$
      gives
      \[
        \frac{3k(k{+}1)^2}{8}
        \le \frac{3}{8}\,n^{3/2} + \frac{15}{8}\,n + 3\sqrt{n} + \frac{3}{2}.
      \]
  \end{enumerate}
  Combining (i)--(iii):
  \begin{equation}\label{eq:lb-allterms}
    a(n) \ge \frac{n^2}{2} - \frac{23}{8}\,n^{3/2}
      - \frac{11}{8}\,n - 3\sqrt{n} - \frac{3}{2}.
  \end{equation}

  \textit{Tail bound.}  We claim
  \begin{equation}\label{eq:tail}
    \frac{11}{8}\,n + 3\sqrt{n} + \frac{3}{2}
    \le \frac{3}{8}\,n^{3/2}
    \qquad (n \ge 36).
  \end{equation}
  Equivalently, $11n + 24\sqrt{n} + 12 \le 3n^{3/2}$.  At $n = 36$
  the right side is $648$ and the left is $552$, so the inequality
  holds.  The derivative of the difference,
  $\frac{9}{2}\sqrt{n} - 11 - 12/\sqrt{n}$, is positive for
  $n \ge 36$ (at $n = 36$: $27 - 11 - 2 = 14 > 0$) and increasing.
  Hence~\eqref{eq:tail} holds for all $n \ge 36$.

  Combining~\eqref{eq:lb-allterms} and~\eqref{eq:tail}:
  \[
    a(n) \ge \frac{n^2}{2} - \frac{23}{8}\,n^{3/2}
    - \frac{3}{8}\,n^{3/2}
    = \frac{n^2}{2} - \frac{13}{4}\,n^{3/2}. \qedhere
  \]
\end{proof}

\begin{corollary}\label{cor:rate}
  For all $n \ge 36$:
  $\;\frac{1}{2} - \frac{13}{4}\,n^{-1/2}
  \le a(n)/n^2
  \le \frac{1}{2} + \frac{1}{2n}$.
\end{corollary}

\begin{remark}\label{rem:optimal-constant}
  The constant $13/4$ is not expected to be optimal.  Optimising the
  construction and its partition bookkeeping is a separate problem; the
  present paper uses only the explicit bound proved in
  Theorem~\ref{thm:quant-sharp}.
\end{remark}

\begin{remark}\label{rem:comparison}
  Theorem~\ref{thm:quant-sharp} is non-trivially positive for
  $n \ge 43$, compared to $n \ge 901$ for Theorem~\ref{thm:quant};
  the improvement is mainly useful because it makes the asymptotic
  estimate visible already in the exact computational range.
\end{remark}

% =======================================================================
\section{Density of ones}\label{sec:density}
% =======================================================================

\subsection{Near-optimal strings}

\begin{proposition}\label{prop:near-opt-density}
  There exists a family of strings $s_n$ with
  $f(s_n) \ge n^2/2 - O(n^{3/2})$ and
  $\ones(s_n)/n = 1 - O(n^{-1/2})$.
\end{proposition}

\begin{proof}
  For $n \ge 64$, take $s_n$ from~\eqref{eq:string}.  By Theorem~\ref{thm:quant},
  $f(s_n) \ge n^2/2 - 15n^{3/2}$.  The string contains exactly
  $k - 1 \le \sqrt{n}/4$ zeros, so
  $\ones(s_n)/n \ge 1 - 1/(4\sqrt{n})$.
  The finitely many smaller values of~$n$ may be filled in arbitrarily.
\end{proof}

\subsection{The zero-position penalty}

\begin{lemma}[Zero-position bound]\label{lem:zero-pos}
  Let $s \in \{0,1\}^n$ with $z \ge 1$ zeros at positions
  $p_1 < \cdots < p_z$.  Then
  \[
    f(s) \le \frac{n(n{+}1)}{2} - \sum_{r=1}^z (n - p_r + 1) + 1.
  \]
  In particular, $\;f(s) \le n(n{+}1)/2 - z(z{+}1)/2 + 1$.
\end{lemma}

\begin{proof}
  Since $z \ge 1$, the value $0$ is in $V(s)$ (any single zero bit
  is a substring representing~$0$).  Every other element of $V(s)$ is
  positive; for each such~$x$, choose a minimum-length representative
  $w_x$.  By Lemma~\ref{lem:canon} (stated in
  Section~\ref{sec:convexity}), $w_x$ starts with~$\mathtt{1}$.
  This map is injective, so distinct positive values correspond to
  distinct leading-$\mathtt{1}$ substrings.

  Each zero at position $p_r$ eliminates $n - p_r + 1$ potential
  leading-$\mathtt{1}$ starts (all substrings beginning there start
  with~$\mathtt{0}$).  Hence the number of
  leading-$\mathtt{1}$ substrings is at most
  $n(n{+}1)/2 - \sum_r(n - p_r + 1)$, and
  \[
    f(s) = |\{0\}| + |\{x \in V(s) : x > 0\}|
    \le 1 + \frac{n(n{+}1)}{2} - \sum_{r=1}^z(n - p_r + 1).
  \]

  The sum is minimized by placing the zeros in the rightmost available
  positions, since each summand decreases as $p_r$ increases.  Thus
  the minimum is achieved by $p_r = n - z + r$, giving
  $\sum_r(z - r + 1) = z(z{+}1)/2$.
\end{proof}

\subsection{Optimal strings have few zeros}

\begin{proposition}\label{prop:optimal-density}
  Every optimal string $s_n^*$ satisfies
  $\zeros(s_n^*) = O(n^{3/4})$, hence
  $\ones(s_n^*)/n = 1 - O(n^{-1/4})$.
\end{proposition}

\begin{proof}
  Let $z_n = \zeros(s_n^*)$.  By Theorem~\ref{thm:lower},
  $a(n) > n = f(\mathtt{1}^n)$ for $n \ge 5$, so an optimal string
  has $z_n \ge 1$ in the range relevant to the asymptotic claim.  For
  $n \ge 36$, Theorem~\ref{thm:quant-sharp} and Lemma~\ref{lem:zero-pos}
  give
  \[
    \frac{n^2}{2} - \frac{13}{4}\,n^{3/2}
    \le a(n)
    \le \frac{n(n{+}1)}{2} - \frac{z_n(z_n{+}1)}{2} + 1.
  \]
  Hence $z_n^2 \le z_n(z_n{+}1) \le (13/2)\,n^{3/2} + n + 4$,
  giving $z_n = O(n^{3/4})$.
\end{proof}

% =======================================================================
\section{Upper bounds from run structure}\label{sec:runs}
% =======================================================================

Lemma~\ref{lem:zero-pos} treats all zeros equally.  We now refine
the bound using the \emph{run structure} of the zeros.

\begin{lemma}[Run-sensitive bound]\label{lem:run-sensitive}
  Let $s \in \{0,1\}^n$ with $z \ge 1$ zeros.  Let the maximal
  zero-runs of $s$ have lengths $c_1,\ldots,c_m$,
  with the $i$-th run beginning at position $a_i$ and
  $r_i = n - (a_i + c_i) + 1$ positions to its right.  Then
  \[
    f(s) \le \frac{n(n{+}1)}{2}
      - \sum_{i=1}^m \Bigl(c_i r_i + \binom{c_i{+}1}{2}\Bigr) + 1.
  \]
\end{lemma}

\begin{proof}
  Apply Lemma~\ref{lem:zero-pos} and group by zero-runs.
  The $i$-th run occupies positions $a_i,\ldots,a_i + c_i - 1$,
  contributing
  $\sum_{t=0}^{c_i-1}(n - a_i - t + 1)
  = c_i(n - a_i + 1) - c_i(c_i{-}1)/2
  = c_i r_i + \binom{c_i{+}1}{2}$
  to the penalty.
\end{proof}

\begin{corollary}[Fragmentation penalty]\label{cor:fragment}
  If $s$ has $z \ge 1$ zeros in $m$ maximal runs, then
  \[
    f(s) \le \frac{n(n{+}1)}{2} - \frac{z(z{+}1)}{2} - \binom{m}{2} + 1.
  \]
\end{corollary}

\begin{proof}
  Since consecutive zero-runs are separated by at least
  one~$\mathtt{1}$, for $i < m$:
  $r_i \ge \sum_{j>i} c_j + (m - i)$.  Hence
  $\sum_i c_i r_i \ge \sum_{i<j} c_i c_j + \sum_{i=1}^{m-1} c_i(m{-}i)$.
  The identity
  $\sum_{i<j} c_i c_j + \sum_i \binom{c_i{+}1}{2} = \binom{z{+}1}{2}$
  holds since $z = \sum c_i$.
  Finally, $\sum_{i=1}^{m-1} c_i(m{-}i) \ge \sum_{i=1}^{m-1}(m{-}i)
  = \binom{m}{2}$ since each $c_i \ge 1$.
\end{proof}

\begin{remark}
  Corollary~\ref{cor:fragment} shows that fragmenting zeros into more
  runs always reduces the upper bound, regardless of positions.
  Each additional zero-run incurs a penalty of at least $m - 1$.
\end{remark}

% =======================================================================
\section{The run-length frequency bound}\label{sec:rlf}
% =======================================================================

\begin{theorem}[Run-length frequency bound]\label{thm:rlf}
  Every $s \in \{0,1\}^n$ satisfies
  \begin{equation}\label{eq:rlf}
    f(s) \le 1 + L(s) + \binom{n{+}1}{2}
      - \sum_{t=1}^p \binom{\ell_t{+}1}{2}
      - \sum_{t=1}^q \binom{m_t{+}1}{2}.
  \end{equation}
\end{theorem}

\begin{proof}
  Let $V_1(s)$ be the set of values representable by an
  all-$\mathtt{1}$ substring, and let $V_{\mathrm{m}}(s)$ be the set
  of remaining positive values.  Then $V(s)$ is contained in the
  disjoint union $\{0\} \cup V_1(s) \cup V_{\mathrm{m}}(s)$; the
  singleton $\{0\}$ is harmless even when $0 \notin V(s)$.

  \textit{Counting $V_1$.}
  An all-$\mathtt{1}$ substring of length~$j$ has value $2^j - 1$,
  depending only on~$j$.  Hence
  $V_1(s) = \{2^j - 1 : 1 \le j \le L(s)\}$ and
  $|V_1(s)| = L(s)$.

  \textit{Counting $V_{\mathrm{m}}$.}
  Every value in $V_{\mathrm{m}}(s)$ is represented by a substring
  containing both~$\mathtt{0}$ and~$\mathtt{1}$.  Hence
  $|V_{\mathrm{m}}(s)|$ is at most the number of pairs $(i,j)$ with
  $1 \le i \le j \le n$ such that $s[i..j]$ contains both symbols.
  This equals $\binom{n{+}1}{2}$ minus the count of monochromatic
  substring pairs.  A maximal $1$-run of length~$\ell_t$ contributes
  $\binom{\ell_t{+}1}{2}$ monochromatic pairs, and similarly for
  $0$-runs.  Hence
  $|V_{\mathrm{m}}| \le \binom{n{+}1}{2}
  - \sum_t \binom{\ell_t{+}1}{2} - \sum_t \binom{m_t{+}1}{2}$.

  Combining via $f(s) \le 1 + |V_1| + |V_{\mathrm{m}}|$
  gives~\eqref{eq:rlf}.
\end{proof}

\begin{corollary}[Concentration of run lengths]\label{cor:concentration}
  For any optimum~$s$ of $a(n)$ with $n \ge 36$:
  \[
    \sum_{t=1}^p \binom{\ell_t{+}1}{2}
    + \sum_{t=1}^q \binom{m_t{+}1}{2}
    \le \frac{13}{4}\,n^{3/2} + \frac{n}{2} + L(s) + 1.
  \]
\end{corollary}

\begin{proof}
  Rearrange Theorem~\ref{thm:rlf} with $f(s) = a(n)$ and apply the
  lower bound $a(n) \ge n^2/2 - (13/4)\,n^{3/2}$
  (Theorem~\ref{thm:quant-sharp}), using
  $\binom{n{+}1}{2} = n^2/2 + n/2$.
\end{proof}

\begin{corollary}[Longest $1$-run]\label{cor:Lbound}
  For every optimum~$s$ of $a(n)$ with $n \ge 36$:
  $\;L(s) = O(n^{3/4})$.
\end{corollary}

\begin{proof}
  The term $\binom{L{+}1}{2}$ appears on the left of
  Corollary~\ref{cor:concentration}; all other terms are non-negative.
  Hence $\binom{L{+}1}{2} \le (13/4)\,n^{3/2} + n/2 + L + 1$,
  giving $L^2/2 \le (13/4)\,n^{3/2} + L + n/2 + 1$ and thus
  $L \le \sqrt{13/2}\,n^{3/4} + O(n^{3/8})$.
\end{proof}

\begin{corollary}[Longest $0$-run]\label{cor:Kbound}
  For every optimum~$s$ of $a(n)$ with $n \ge 36$:
  $\;K(s) = O(n^{3/4})$.
\end{corollary}

\begin{proof}
  The term $\binom{K{+}1}{2}$ appears on the left of
  Corollary~\ref{cor:concentration}.  By Corollary~\ref{cor:Lbound},
  $L(s) = O(n^{3/4})$, so
  \[
    \binom{K{+}1}{2}
    \le \frac{13}{4}\,n^{3/2} + \frac{n}{2} + O(n^{3/4}) + 1.
  \]
  Hence $K^2 = O(n^{3/2})$, giving $K = O(n^{3/4})$.
\end{proof}

% =======================================================================
\section{Convexity of the sequence}\label{sec:convexity}
% =======================================================================

\begin{theorem}[Convexity]\label{thm:convex}
  For every $n \ge 2$,
  $\;a(n{+}1) - 2a(n) + a(n{-}1) \ge 0$.
  Equivalently, $\Da(n)$ is nondecreasing.
\end{theorem}

The proof requires two preliminary lemmas.

\begin{lemma}[Canonical witness]\label{lem:canon}
  For any $s \in \{0,1\}^n$ and any positive $x \in V(s)$, the
  standard binary representation of~$x$ (the unique string starting
  with~$\mathtt{1}$ of length $\lfloor \log_2 x \rfloor + 1$)
  occurs as a substring of~$s$.  Moreover, every minimum-length
  substring of~$s$ representing~$x$ equals this standard
  representation.
\end{lemma}

\begin{proof}
  Let $w_x$ be any minimum-length substring of~$s$ representing~$x$.
  If $w_x$ started with~$\mathtt{0}$, removing the leading zero
  gives a shorter representative, contradicting minimality.  So
  $w_x$ starts with~$\mathtt{1}$ and has length at least
  $\lfloor \log_2 x \rfloor + 1$.  By minimality, equality holds.
  Since both $w_x$ and the standard representation start
  with~$\mathtt{1}$, have the same length, and represent~$x$, they
  are identical.  In particular, the standard representation occurs
  as a substring of~$s$.
\end{proof}

\begin{lemma}[Zero dichotomy]\label{lem:zero-dichotomy}
  Let $s \in \{0,1\}^n$ with $n \ge 2$, and let
  $p = s_1 \cdots s_{n-1}$, $q = s_2 \cdots s_n$.
  Then $0 \notin V(s) \setminus V(p)$ or
  $0 \notin V(s) \setminus V(q)$ (or both).
\end{lemma}

\begin{proof}
  If $0 \in V(s) \setminus V(p)$, then $s$ contains a zero but $p$
  does not, so the only zero of~$s$ is at position~$n$.
  If $0 \in V(s) \setminus V(q)$, then $s$ contains a zero but $q$
  does not, so the only zero of~$s$ is at position~$1$.
  Since $n \ge 2$, positions $1$ and $n$ are distinct, so both
  conditions cannot hold simultaneously.
\end{proof}

\begin{proof}[Proof of Theorem~\ref{thm:convex}]
  Fix $n \ge 2$ and let $s = s_1 \cdots s_n$ be optimal for~$a(n)$.
  Let $p = s_1 \cdots s_{n-1}$ and $q = s_2 \cdots s_n$.  Since
  $V(p) \subseteq V(s)$ and $V(q) \subseteq V(s)$:
  \begin{equation}\label{eq:Vdiff}
    |V(s) \setminus V(p)| \ge a(n) - a(n{-}1) = \Da(n), \qquad
    |V(s) \setminus V(q)| \ge a(n) - a(n{-}1) = \Da(n).
  \end{equation}

  By Lemma~\ref{lem:zero-dichotomy}, at least one of
  $V(s) \setminus V(p)$ and $V(s) \setminus V(q)$ contains only
  positive values.  We show $a(n{+}1) \ge a(n) + \Da(n)$ in each
  case.

  \medskip
  \noindent\textbf{Case 1: $0 \notin V(s) \setminus V(p)$.}
  Every $x \in V(s) \setminus V(p)$ is positive.  By
  Lemma~\ref{lem:canon}, the standard binary representation~$w_x$
  (of length $\ell_x = \lfloor \log_2 x \rfloor + 1$, starting
  with~$\mathtt{1}$) occurs as a substring of~$s$.

  \textit{$w_x$ must end at position~$n$.}
  If $w_x$ ended at position $j \le n - 1$, it would lie entirely
  within~$p$, giving $x \in V(p)$, a contradiction.

  Form $t = s\,\mathtt{0}$ of length $n + 1$.  For each
  $x \in V(s) \setminus V(p)$, the string $w_x\mathtt{0}$ is a
  substring of~$t$.

  \textit{$w_x\mathtt{0}$ is the standard representation of $2x$.}
  Since $w_x$ starts with~$\mathtt{1}$ and
  $2^{\ell_x - 1} \le x < 2^{\ell_x}$, we have
  $2^{\ell_x} \le 2x < 2^{\ell_x + 1}$, so $w_x\mathtt{0}$ starts
  with~$\mathtt{1}$ and has length
  $\lfloor \log_2(2x) \rfloor + 1 = \ell_x + 1$.

  \textit{$2x \notin V(s)$.}  Suppose $2x \in V(s)$.  By
  Lemma~\ref{lem:canon}, $s$ contains the standard representation
  $w_x\mathtt{0}$ as a substring, say occupying positions
  $i, \ldots, i + \ell_x$.  Then $w_x$ occupies positions
  $i, \ldots, i + \ell_x - 1$ with $i + \ell_x \le n$, so
  $w_x$ ends at position $i + \ell_x - 1 \le n - 1$, placing $w_x$
  inside~$p$.  This gives $x \in V(p)$, a contradiction.

  \textit{Injectivity.}  The map $x \mapsto 2x$ is injective.

  \textit{Conclusion.}  The values
  $\{2x : x \in V(s) \setminus V(p)\}$ are pairwise distinct,
  lie in $V(t) \setminus V(s)$, and number at least~$\Da(n)$.
  Hence $f(t) \ge a(n) + \Da(n)$.

  \medskip
  \noindent\textbf{Case 2: $0 \notin V(s) \setminus V(q)$.}
  Every $x \in V(s) \setminus V(q)$ is positive.  By
  Lemma~\ref{lem:canon}, the standard representation~$w_x$ occurs
  in~$s$.

  \textit{$w_x$ must start at position~$1$.}
  If $w_x$ started at position $i \ge 2$, it would lie entirely
  within $q = s_2 \cdots s_n$, giving $x \in V(q)$, a contradiction.

  Form $t = \mathtt{1}\,s$ of length $n + 1$.  For each
  $x \in V(s) \setminus V(q)$, the string $\mathtt{1}\,w_x$ is a
  substring of~$t$ (occupying positions $1, \ldots, \ell_x + 1$).
  Set $y = 2^{\ell_x} + x$.

  \textit{$\mathtt{1}\,w_x$ is the standard representation of~$y$.}
  Since $w_x$ starts with~$\mathtt{1}$, the string
  $\mathtt{1}\,w_x$ starts with~$\mathtt{1}$ and has length
  $\ell_x + 1$.  Now
  $2^{\ell_x} \le y < 2^{\ell_x + 1}$, so the length is correct.

  \textit{$y \notin V(s)$.}  If $y \in V(s)$, by
  Lemma~\ref{lem:canon} the string $\mathtt{1}\,w_x$ occurs in~$s$,
  say starting at position~$i$.  Then $w_x$ occupies positions
  $i+1, \ldots, i + \ell_x$ with $i + 1 \ge 2$, so $w_x$ lies
  within $q = s_2 \cdots s_n$, giving $x \in V(q)$, a contradiction.

  \textit{Injectivity.}  If $\ell_x = \ell_{x'}$ then
  $y \ne y'$ since $x \ne x'$.  If $\ell_x < \ell_{x'}$ then
  $y < 2^{\ell_x + 1} \le 2^{\ell_{x'}} \le y'$.

  \textit{Conclusion.}  The values
  $\{2^{\ell_x} + x : x \in V(s) \setminus V(q)\}$ are pairwise
  distinct, lie in $V(t) \setminus V(s)$, and number at
  least~$\Da(n)$.
  Hence $f(t) \ge a(n) + \Da(n)$.

  \medskip
  In both cases,
  $a(n{+}1) \ge a(n) + \Da(n) = 2a(n) - a(n{-}1)$.
\end{proof}

\begin{remark}\label{rem:convex-upper}
  Theorem~\ref{thm:convex} establishes $\DDa(n) \ge 0$.
  Computational data further suggest $\DDa(n) \le 2$ for all
  $n \ge 2$ (see Section~\ref{sec:observations}).
\end{remark}

We now derive a non-trivial upper bound on $\DDa(n)$ by combining
convexity with the global asymptotic lower bound.

\begin{lemma}[First-difference upper bound]\label{lem:da-upper}
  For all $n \ge 1$, $\;\Da(n) \le n$.
\end{lemma}

\begin{proof}
  Let $t$ be any string of length~$n$ and $p = t_1 \cdots t_{n-1}$.
  Every value in $V(t) \setminus V(p)$ is represented only by
  substrings ending at position~$n$, of which there are at most~$n$.
  Hence $f(t) \le f(p) + n \le a(n{-}1) + n$, so
  $a(n) \le a(n{-}1) + n$.
\end{proof}

\begin{theorem}[Second-difference upper bound]\label{thm:dda-upper}
  Define the \emph{deficit} $E(n) = n(n{+}1)/2 - a(n)$.  Then for
  all $n \ge 2$ and every integer $1 \le h < n$,
  \[
    \DDa(n) \le \frac{h{+}1}{2} + \frac{E(n)}{h}.
  \]
  Optimising in $h$ gives
  $\;\DDa(n) \le \sqrt{2E(n)} + O(1)$.
\end{theorem}

\begin{proof}
  By convexity (Theorem~\ref{thm:convex}), $\Da$ is nondecreasing,
  so for any $h \ge 1$:
  \[
    a(n) - a(n{-}h) = \sum_{j=0}^{h-1} \Da(n{-}j) \le h\,\Da(n).
  \]
  Since $a(n{-}h) \le (n{-}h)(n{-}h{+}1)/2$, we obtain
  \begin{align*}
    \Da(n) &\ge \frac{a(n) - a(n{-}h)}{h}
    \ge \frac{n(n{+}1)/2 - E(n) - (n{-}h)(n{-}h{+}1)/2}{h} \\
    &= \frac{h(2n{-}h{+}1)/2 - E(n)}{h}
    = n - \frac{h{-}1}{2} - \frac{E(n)}{h}.
  \end{align*}
  Combining with Lemma~\ref{lem:da-upper}:
  \[
    \DDa(n) = \Da(n{+}1) - \Da(n)
    \le (n{+}1) - \Bigl(n - \frac{h{-}1}{2} - \frac{E(n)}{h}\Bigr)
    = \frac{h{+}1}{2} + \frac{E(n)}{h}.
  \]
  Setting $h = \lfloor \sqrt{2E(n)} \rfloor$ gives
  $\DDa(n) \le \sqrt{2E(n)} + O(1)$.
\end{proof}

\begin{corollary}\label{cor:dda-sublinear}
  $\DDa(n) = O(n^{3/4})$.  More precisely, for $n \ge 36$,
  \[
    \DDa(n) \le \sqrt{\tfrac{13}{2}\,n^{3/2} + n} + O(1).
  \]
\end{corollary}

\begin{proof}
  By Theorem~\ref{thm:quant-sharp},
  $E(n) \le \tfrac{13}{4}\,n^{3/2} + \tfrac{n}{2}$ for $n \ge 36$.
  Apply Theorem~\ref{thm:dda-upper}.
\end{proof}

\begin{remark}
  Any improvement in the global deficit bound
  $E(n) = n(n{+}1)/2 - a(n)$ immediately tightens the second-difference
  bound via Theorem~\ref{thm:dda-upper}.  If the conjectured
  $E(n) = \Theta(n^{3/2})$ is optimal, the method yields
  $\DDa(n) = O(n^{3/4})$, far short of the empirical bound
  $\DDa(n) \le 2$.  Closing that gap appears to require additional
  structural information.
\end{remark}

% =======================================================================
\section{The ones-weighted crossing bound}\label{sec:crossing}
% =======================================================================

The following bound, which partitions the score of a string into a
prefix contribution and a residual, provides the theoretical
foundation for the branch-and-bound exact solver described in
Section~\ref{sec:computation}.

\begin{lemma}[Ones-weighted crossing bound]\label{lem:crossing}
  Let $p \in \{0,1\}^k$ be a prefix and
  $e \in \{0,1\}^m$ any extension, with $n = k + m$ and $s = pe$.
  Let $V_0(p) = \{\val(p[i..j]) : 1 \le i \le j \le k\}$.  Then
  \[
    f(s) \le |V_0(p)| + a(m) + \ones(p) \cdot m,
  \]
  where $\ones(p) = \#\{i : p_i = 1\}$.
\end{lemma}

\begin{proof}
  Write $V_+(p,e)$ for the set of values represented by substrings
  not entirely within~$p$.  Then
  $f(s) \le |V_0(p)| + |V_+(p,e)|$.

  Partition $V_+(p,e)$ by where each defining substring begins:

  \textit{(i) Suffix-only substrings} (starting at position $> k$)
  contribute at most $|V(e)| \le a(m)$ values.

  \textit{(ii) Zero-leading crossing substrings} (starting at $i \le k$
  with $p_i = 0$, ending at $j > k$):  stripping leading zeros either
  reduces to a suffix-only value or to a one-leading crossing value.
  These contribute no new values beyond those in (i) and (iii).

  \textit{(iii) One-leading crossing substrings} (starting at $i \le k$
  with $p_i = 1$, extending $j \in \{1,\ldots,m\}$ bits into~$e$):
  for fixed~$i$, set $A_i = \val(p[i..k]) \ge 1$.  The value is
  $v_{i,j} = A_i \cdot 2^j + \val(e[1..j])$.  For $j_1 < j_2$:
  $v_{i,j_2} \ge A_i \cdot 2^{j_2} \ge 2A_i \cdot 2^{j_1}
  > (A_i + 1) \cdot 2^{j_1} - 1 \ge v_{i,j_1}$
  (using $A_i \ge 1$), so the $m$ values are strictly increasing in~$j$.
  There are $\ones(p)$ one-positions, giving at most
  $\ones(p) \cdot m$ crossing values.

  Combining: $|V_+| \le a(m) + \ones(p) \cdot m$.
\end{proof}

\begin{remark}
  In the branch-and-bound solver, this bound is evaluated at every
  node of the search tree: the prefix~$p$ is the partially
  constructed string, and $a(m)$ is looked up from previously proven
  values or from Theorem~\ref{thm:length-strat}.  The bound is
  typically $30$--$50\%$ tighter than the naive residual
  $n(n{+}1)/2 - |V_0(p)|$.
\end{remark}

\subsection{Suffix-automaton refinement}

The crossing bound of Lemma~\ref{lem:crossing} counts $\ones(p) \cdot m$
crossing values as potentially new.  Many of these are already present
inside the fixed prefix and can be detected via the suffix automaton
(SAM) of~$p$.

\begin{lemma}[SAM savings]\label{lem:sam-savings}
  Let $p \in \{0,1\}^{k}$, let $e \in \{0,1\}^m$, put $s=pe$, and let
  $\mathcal{S}$ be the suffix automaton of~$p$.  For each position
  $i \le k$ with $p_i = 1$, define $d_i$ to be the largest
  integer $d$ such that the SAM state reached by the suffix $p[i..k]$
  has a complete binary transition subtree of depth~$d$.  Equivalently,
  for every $1 \le j \le d$ and every binary string
  $w \in \{0,1\}^j$, the word $p[i..k] \cdot w$ occurs in~$p$, so its
  value already lies in $V_0(p)$.  Then
  \[
    f(s) \le |V_0(p)| + a(m) + \sum_{\substack{i=1 \\ p_i=1}}^{k}
    (m - d_i)^+,
  \]
  where $(x)^+ = \max(x, 0)$.
\end{lemma}

\begin{proof}
  By the proof of Lemma~\ref{lem:crossing}, each one-position
  $i$ contributes at most $m$ crossing values $v_{i,1},\ldots,v_{i,m}$.
  The suffix automaton represents exactly the substrings of the fixed
  prefix.  Therefore, if the state reached by the suffix $p[i..k]$ has
  all binary paths of length~$j$, then for every $w \in \{0,1\}^j$ the
  word $p[i..k]w$ occurs as a substring of the
  fixed prefix and its value already lies in $V_0(p)$.  For
  $j = 1,\ldots,d_i$, these are precisely the first $d_i$ possible
  crossing extensions from position~$i$, so they are guaranteed
  duplicates.  Subtracting them from the budget of $m$ possible new
  crossing values gives at most $(m-d_i)^+$ new values from that
  position.  Summing over all one-positions yields the result.
\end{proof}

\begin{remark}
  Computing all $d_i$ via a bounded breadth-first search of depth
  $J \le 3$ in the SAM takes $O(k \cdot 2^J)$ time per node of the
  search tree.  In practice, this refinement reduces the crossing
  budget by $20$--$40\%$ at moderate prefix lengths, and is the
  primary reason the branch-and-bound solver terminates within
  practical time for $n$ up to~$111$.
\end{remark}

% =======================================================================
\section{Separator structure}\label{sec:separator}
% =======================================================================

The certified optimal strings for $81 \le n \le 111$ all have
separator sequence beginning $(1,2,1,1,\ldots)$.  We prove below that
any two consecutive separators of equal length produce forced
collisions (Theorem~\ref{thm:sep2-collisions}).  This gives a rigorous
obstruction to repeated adjacent separator lengths, but it is not by
itself a proof of the full separator pattern.

\subsection{Forced collisions from equal consecutive separators}

The following theorem shows that two consecutive separators of the
same length force a quadratic number of collisions.

\begin{theorem}[Forced collisions from equal consecutive separators]
\label{thm:sep2-collisions}
  Let $w$ be a binary string with ones-block lengths
  $b_1, b_2, b_3, \ldots$ and separator lengths $s_1, s_2, \ldots$,
  so that $w$ begins
  $\mathtt{1}^{b_1}\,\mathtt{0}^{s_1}\,\mathtt{1}^{b_2}\,
   \mathtt{0}^{s_2}\,\mathtt{1}^{b_3}\,\cdots$.
  Suppose two consecutive separators have the same length:
  $s_i = s_{i+1} = s$ for some~$i$.  Set
  $J = \min(b_i, b_{i+1})$ and $K = \min(b_{i+1}, b_{i+2})$.  Then the
  $J \cdot K$ values
  \[
    \bigl\{v_{j,k} = \val(\mathtt{1}^j\,\mathtt{0}^s\,\mathtt{1}^k) :
    1 \le j \le J,\; 1 \le k \le K\bigr\}
  \]
  are pairwise distinct and each appears at least twice as a
  substring of~$w$.  Explicitly,
  $v_{j,k} = (2^j - 1)\cdot 2^{s+k} + 2^k - 1 = 2^{j+s+k} - 2^{s+k} + 2^k - 1$.
\end{theorem}

\begin{proof}
  \textbf{Two occurrences.}
  For $j \le J$ and $k \le K$, the pattern
  $P_{j,k} = \mathtt{1}^j\,\mathtt{0}^s\,\mathtt{1}^k$ (representing
  $v_{j,k}$) appears at two positions in~$w$:
  \begin{enumerate}[nosep]
    \item Crossing separator~$i$: the last $j$ ones of
      block~$i$ ($j \le b_i$), then $\mathtt{0}^s = \mathtt{0}^{s_i}$,
      then the first $k$ ones of block~$i{+}1$ ($k \le b_{i+1}$).
    \item Crossing separator~$i{+}1$: the last
      $j$ ones of block~$i{+}1$ ($j \le b_{i+1}$), then
      $\mathtt{0}^s = \mathtt{0}^{s_{i+1}}$, then the first
      $k$ ones of block~$i{+}2$ ($k \le b_{i+2}$).
  \end{enumerate}
  These positions differ since $b_{i+1} \ge 1$.

  \textbf{Distinctness.}
  Each $v_{j,k}$ has standard binary representation
  $\mathtt{1}^j\mathtt{0}^s\mathtt{1}^k$ (a string of length
  $j + s + k$ starting with~$\mathtt{1}$).  If
  $v_{j_1,k_1} = v_{j_2,k_2}$, their standard representations are
  identical, so $j_1 + k_1 = j_2 + k_2$ (same total length) and the
  $\mathtt{0}^s$ block occupies the same positions: $j_1 = j_2$,
  hence $k_1 = k_2$.

  \textbf{These are genuine collisions.}
  Each $v_{j,k}$ appears at two distinct starting positions, so its
  substring occurrence is counted twice but contributes only once to
  $f(w)$.  This represents $J \cdot K$ forced collisions.
\end{proof}

\begin{corollary}\label{cor:sep2-deficit}
  If a string $w$ with $z \ge 1$ zeros has $s_i = s_{i+1}$ for some~$i$,
  then
  $f(w) \le n(n{+}1)/2 - Z(w) + 1 - C(w) - J \cdot K$,
  where $p_1<\cdots<p_z$ are the zero positions,
  $Z(w) = \sum_{r=1}^z(n - p_r + 1)$ is the zero-position penalty
  (Lemma~\ref{lem:zero-pos}),
  $C(w) = \sum_j \binom{b_j+1}{2} - \max_j b_j$ counts one-block
  self-collisions, and $J, K$ are as in
  Theorem~\ref{thm:sep2-collisions}.
\end{corollary}

\begin{proof}
  Let $N = n(n{+}1)/2$.  By Lemma~\ref{lem:canon}, every positive
  value has a standard representative beginning with~$\mathtt{1}$.
  There are at most $N - Z(w)$ leading-$\mathtt{1}$ substrings, and
  the value~$0$ contributes at most one further value.

  Among the leading-$\mathtt{1}$ substrings, the all-$\mathtt{1}$
  substrings contribute $\sum_j \binom{b_j+1}{2}$ occurrences but at
  most $\max_j b_j$ distinct values, so they account for at least
  $C(w)$ duplicate occurrences.

  The $J \cdot K$ collisions of Theorem~\ref{thm:sep2-collisions} are
  disjoint from these all-$\mathtt{1}$ collisions because their
  standard representatives
  $P_{j,k} = \mathtt{1}^j\,\mathtt{0}^s\,\mathtt{1}^k$ contain a
  zero.  They are also disjoint from zero-starting substrings because
  each $P_{j,k}$ begins with~$\mathtt{1}$.  Therefore the leading-one
  substring count overcounts the positive values by at least
  $C(w) + J K$, and
  \[
    f(w) \le 1 + (N - Z(w)) - C(w) - J K. \qedhere
  \]
\end{proof}

\begin{remark}
  The theorem requires only that $s_i = s_{i+1}$; it applies to
  consecutive separators of any common length.
  Making two consecutive separators differ eliminates the collisions:
  the two occurrences of $P_{j,k}$ then cross zero-blocks of different
  lengths, producing the distinct values
  $\val(\mathtt{1}^j\,\mathtt{0}^{s_i}\,\mathtt{1}^k) \ne
   \val(\mathtt{1}^j\,\mathtt{0}^{s_{i+1}}\,\mathtt{1}^k)$.

  For $n \approx 100$ with hypothetical separator sequence
  $(1, 1, \ldots)$, the first pair $(s_1, s_2) = (1, 1)$ gives
  $J \cdot K = \min(b_1, b_2) \cdot \min(b_2, b_3)
   \approx 17 \cdot 14 = 238$ forced collisions, roughly $20\%$ of
  the total deficit.
  Changing $s_2$ from~$1$ to~$2$ eliminates these collisions at the
  cost of one additional bit; separators $\ge 3$ also work but
  waste an extra position.  In the computed optima, later separator
  pairs have adjacent blocks small enough that the collision count
  $\min(b_i, b_{i+1}) \cdot \min(b_{i+1}, b_{i+2})$ is negligible,
  making the cost of an extra separator bit outweigh the observed
  savings.  This is consistent with the pattern
  $(1, 2, 1, 1, 1, \ldots)$ favoured by the data, but the infinite
  statement remains
  Conjecture~\ref{conj:seps}.
\end{remark}

% =======================================================================
\section{Computational results}\label{sec:computation}
% =======================================================================

\subsection{Scoring}\label{sec:scoring}

For small $n$ ($n \lesssim 30$), $f(s)$ is computed by inserting all
substring values into a hash set.  For larger~$n$, we use the suffix
array~$\mathrm{SA}$ and LCP array (via Kasai's
algorithm~\cite{kasai2001}):
\[
  f(s) = \sum_{\substack{0 \le i < n \\ s[\mathrm{SA}[i]] = 1}}
    (n - \mathrm{SA}[i] - \mathrm{LCP}[i])
    + \mathbf{1}[\mathtt{0} \in s].
\]
Here the suffix array is indexed from~$0$ and $\mathrm{LCP}[0]=0$.
This formula counts, for each suffix in sorted order, the number of
new leading-$\mathtt{1}$ substrings it contributes beyond those
shared with the preceding suffix (as measured by the LCP), then adds
$1$ if the string contains any zero (for the value~$0$).  The
standard suffix-array argument applies because every positive value is
identified with its unique leading-$\mathtt{1}$ representation:
prefixes of a suffix that are not longer than its LCP with the
preceding suffix have already occurred, and longer prefixes are new.
Thus one scoring pass for a fixed string takes $O(n \log n)$ time, or
$O(n)$ when the suffix array is constructed with the DivSufSort
library~\cite{pydivsufsort}.  This is distinct from the size of the
branch-and-bound search tree considered below.

\subsection{Branch-and-bound exact solver}\label{sec:bb-solver}

To prove exact values of $a(n)$ beyond $n = 80$, we developed a
branch-and-bound solver that constructs $n$-bit strings bit-by-bit,
pruning branches whose upper bound falls below the current best
known score.

\begin{definition}[Upper bound at a search node]
  At a node where the prefix $p \in \{0,1\}^k$ has been fixed
  (with $m = n - k$ remaining bits), let $B(m)$ be any certified upper
  bound for $a(m)$.  The node upper bound is
  \[
    U(p) = |V_0(p)| + B(m) + \sum_{\substack{i=1 \\ p_i=1}}^{k}
    (m - d_i)^+,
  \]
  where $V_0(p)$ is the set of distinct values of substrings
  entirely within~$p$.  In a self-contained certificate one may take
  $B(m)$ from Theorem~\ref{thm:length-strat}; tighter runs may replace
  this by exact values only when those values are independently
  certified.  The integer $d_i$ is the SAM savings depth at
  position~$i$ (Lemma~\ref{lem:sam-savings}).
\end{definition}

Since $B(m) \ge a(m)$, Lemmas~\ref{lem:crossing} and
\ref{lem:sam-savings} imply that $U(p)$ is a valid upper bound on
$f(s)$ for any completion $s$ of~$p$.

The algorithm (pseudocode below) proceeds as follows:
\begin{enumerate}[nosep]
  \item Initialise the incumbent $L \gets$ a known lower bound (from
    a heuristic search or a previously computed value for a
    nearby~$n$).
  \item Perform depth-first search over all possible bit assignments,
    maintaining a suffix automaton (SAM) for the current prefix to
    track $|V_0(p)|$ and the SAM savings~$d_i$ incrementally.
  \item At each node, compute $U(p)$.  In value-only mode, prune when
    $U(p) \le L$; in \texttt{collect\_all} mode, prune only when
    $U(p) < L$, so branches that may contain further optima are kept.
  \item At a leaf ($k = n$), evaluate $f(s)$ exactly.  If $f(s)>L$,
    update $L$ and, in \texttt{collect\_all} mode, clear the current
    optimum list; if $f(s)=L$, record~$s$ in that list.
  \item On termination, $L = a(n)$.
\end{enumerate}

\begin{center}
\begin{tabular}{l}
\toprule
\textsc{BranchAndBound}$(n)$ \\
\midrule
\textbf{if} $n = 1$: \quad \textbf{return} $1$ \\
$L \gets \text{initial lower bound}$ \\
\textsc{Search}$(\mathtt{1}, \text{SAM}_{\mathtt{1}}, 1, 1)$ \\
\textbf{return} $L$ \\[4pt]
\textsc{Search}$(p,\; \mathcal{S},\; |V_0|,\; \ones(p))$ \\
\quad $k \gets |p|$; \quad $m \gets n - k$ \\
\quad \textbf{if} $k = n$: \quad update $L$ and optimum records; \quad \textbf{return} \\
\quad Compute $d_i$ for each one-position $i$ via SAM~$\mathcal{S}$ \\
\quad $U \gets |V_0| + B(m) + \sum_{i:\, p_i=1} (m - d_i)^+$ \\
\quad \textbf{if} \textsc{Prunable}$(U,L)$: \quad \textbf{return} \\
\quad \textbf{for} $b \in \{1, 0\}$: \\
\quad\quad Extend SAM: $\mathcal{S}' \gets \textsc{SamExtend}(\mathcal{S}, b)$ \\
\quad\quad $|V_0'| \gets$ updated distinct substring count \\
\quad\quad \textsc{Search}$(p \cdot b,\; \mathcal{S}',\; |V_0'|,\; \ones(p) + b)$ \\
\bottomrule
\end{tabular}
\end{center}
Here \textsc{Prunable}$(U,L)$ means $U < L$ in \texttt{collect\_all}
mode and $U \le L$ in value-only mode.

\begin{remark}[Correctness]
  The solver's correctness rests on two properties: (1)~the upper
  bound $U(p) \ge f(s)$ for every completion $s$ of prefix~$p$
  (Lemmas~\ref{lem:crossing} and~\ref{lem:sam-savings}), so pruning
  never discards a global optimum; and (2)~the search exhaustively
  explores all branches not pruned.  The SAM provides exact tracking
  of $|V_0(p)|$ (each SAM state corresponds to an equivalence class
  of right extensions; the total count of distinct substrings is
  maintained incrementally as each bit is appended).

  The value search is run with initial prefix~$\mathtt{1}$, which is
  valid by Lemma~\ref{lem:leading-bit}.  For results that assert
  properties of all optimal strings, the solver is run in
  \texttt{collect\_all} mode: branches with $U(p)=L$ are not pruned
  merely because one incumbent optimum is already known.  Thus every
  completion with score $a(n)$ is either reached and recorded or lies
  in a subtree whose certified upper bound is below~$a(n)$.
\end{remark}

\begin{remark}[Implementation and reproducibility]
  The implementation uses Python~3.13 with Numba~\cite{numba}
  for JIT compilation of the inner scoring loop and the SAM
  operations.  Parallel search is performed across multiple workers,
  each exploring a subtree rooted at a distinct prefix of length
  $\approx 20$.  Workers share the incumbent~$L$ via atomic updates.
  This affects only the order in which leaves are visited; in
  \texttt{collect\_all} mode the certificate is the final unordered set
  of strings attaining the maximum, not the discovery order.
  Typical running times range from seconds ($n \le 90$) to
  $\sim\!48$ hours ($n = 111$) on a 16-core workstation.

  Source code is available at~\cite{github}.  To reproduce the new
  extension, run the branch-and-bound solver for each target~$n$ with
  \texttt{collect\_all=True}; the solver outputs the exact value, the
  complete list of optimal strings, and the number of nodes explored
  and pruned.  The archived outputs for the extension are included in
  the repository under \texttt{results/branch\_bound\_exact/}, with one
  JSON certificate file \texttt{n\_NNNN\_results.json} per value of~$n$;
  these files record the optimal-string list, \texttt{num\_optimal}, and
  the run metadata.  The consolidated source-data file
  \texttt{data/cached\_results.json} records the exact entries through
  $n=111$, including Fuller's $n\le80$ data and the branch-and-bound
  extension.  The companion sequences A112510, A112511, and A156025 are
  then obtained respectively as the minimum, maximum, and cardinality of
  this list; A156022--A156024 follow from the corresponding value counts.
\end{remark}

\subsection{Exact values}

\begin{theorem}[Certified exact computation]\label{thm:certified-computation}
  The values of $a(n)$ are certified for $1 \le n \le 111$.  For each
  such~$n$, the complete list of optimal strings is certified, and
  hence so are the corresponding values of A112510, A112511, and
  A156025.  The values of A156022--A156024 then follow from the
  definitions.
\end{theorem}

\begin{proof}
  Fuller's A$^*$ search certifies the range $1 \le n \le 80$, with the
  corresponding optimal-string and companion-sequence data recorded in
  the source table.  The branch-and-bound extension is then certified
  sequentially for $n=81,82,\ldots,111$.  At the run for a given~$n$,
  every residual length $m<n$ has already been certified, so the
  computation described in Section~\ref{sec:bb-solver}, run in
  \texttt{collect\_all} mode, may use the certified values for smaller
  residuals as exact upper bounds $B(m)=a(m)$.  Hence the upper bound
  $U(p)$ is certified at every
  search node.  Exhaustive exploration of all branches not pruned then
  proves the optimum.  Since \texttt{collect\_all} mode does not prune
  branches with $U(p)=a(n)$, every string attaining the optimum is
  recorded in the corresponding archived certificate in
  \texttt{results/branch\_bound\_exact/}.  Taking the minimum, maximum,
  and cardinality of this list gives A112510, A112511, and A156025,
  respectively.
\end{proof}

The full machine-readable table is included with the source data; the
main text records only representative values and the growth ratio.
\begin{center}
\begin{tabular}{rrrl}
  \toprule
  $n$ & $a(n)$ & $a(n)/n^2$ & source \\
  \midrule
  50  & 899        & 0.3596 & Fuller/OEIS \\
  80  & 2{,}423    & 0.3786 & Fuller/OEIS \\
  100 & 3{,}875    & 0.3875 & certified here \\
  111 & 4{,}826    & 0.3916 & certified here \\
  \bottomrule
\end{tabular}
\end{center}

\subsection{Exploratory lower bounds for \texorpdfstring{$112 \le n \le 200$}{112 <= n <= 200}}

For $112 \le n \le 200$, lower bounds on $a(n)$ were computed using
two independent Metropolis--Hastings (MH) search methods:
\begin{enumerate}[nosep]
  \item \emph{Unconstrained MH} --- random bit-flip Markov chains on
    the full space of $n$-bit strings, accepting moves that do not
    decrease $f(s)$ ($10^5$ restarts, $10^5$ iterations each).
  \item \emph{Constrained MH} --- chains restricted to a fixed
    run-length template derived from empirical block-range
    constraints, with $2 \times 10^5$ steps per chain across $8$
    parallel chains.
\end{enumerate}
Both methods agreed on the best value for every~$n$ in this range.
These results are certified lower bounds (since $f(s)$ is evaluated
exactly), but are not proved to be exact values of~$a(n)$; the
agreement across independent methods is recorded as evidence only and
is not used in any proof.

\subsection{Large-scale certified lower bounds}

For $n \ge 3{,}000$, a suffix-array-guided local search
iteratively improves a single incumbent string by identifying and
resolving the highest-LCP collisions via targeted single-bit swaps.
The table reports exact scores of explicit strings.  Since each score
is computed exactly, every listed value is a rigorous lower bound
on~$a(n)$, not a certified exact value of~$a(n)$.

\begin{center}
\begin{tabular}{rrl}
  \toprule
  $n$ & Lower bound on $a(n)$ & $a(n)/n^2 \ge$ \\
  \midrule
  $3{,}000$            & $4{,}261{,}266$                          & $0.4735$ \\
  $10{,}000$           & $48{,}490{,}001$                         & $0.4849$ \\
  $10^6$               & $498{,}396{,}662{,}200$                  & $0.4984$ \\
  $10^9$               & $499{,}943{,}911{,}277{,}037{,}650$      & $0.49994$ \\
  $2 \times 10^9$      & $1{,}999{,}841{,}243{,}789{,}594{,}574$  & $0.49996$ \\
  \bottomrule
\end{tabular}
\end{center}

The computation at $n = 2 \times 10^9$ scores a two-billion-bit
string ($\approx 238$~MB as a NumPy array) using the DivSufSort
suffix-array library~\cite{pydivsufsort}.

\begin{remark}
  Theorem~\ref{thm:quant-sharp} gives
  $a(n)/n^2 \ge 1/2 - 13/(4\sqrt{n})$, which reaches $0.4875$ only at
  $n = 14{,}400$.  At moderate values of~$n$, the computational lower
  bounds lie well above this conservative theoretical estimate,
  indicating that the constant $13/4$ is not close to the behaviour of
  the explicit strings found in practice.
\end{remark}

% =======================================================================
\section{Empirical observations and open problems}\label{sec:observations}
% =======================================================================

\subsection{Second differences}

\begin{observation}\label{obs:dda}
  In the certified data through $n \le 111$:
  $\DDa(n) \in \{0, 1, 2\}$.  The value~$2$ occurs at
  $n = 52$ and $n = 71$.  For heuristic values up to $n = 200$,
  $\DDa(n) \in \{0, 1, 2\}$ continues to hold, with additional
  occurrences of~$2$ at $n \in \{117, 157, 178, 191, 194\}$.
\end{observation}

Theorem~\ref{thm:convex} establishes $\DDa(n) \ge 0$ for $n \ge 2$.

\begin{conjecture}\label{conj:dda}
  $\DDa(n) \le 2$ for all $n \ge 2$.
\end{conjecture}

\subsection{The separator pattern}

\begin{observation}\label{obs:seps}
  The certified data show that, for every $8 \le n \le 111$, every
  optimal $n$-bit string has first separator length~$1$.
  The threshold is sharp: exhaustive enumeration for $n \le 7$ shows
  that, for $n = 4,5,6,7$, some optimal strings have first separator
  length greater than~$1$.

  For all $81 \le n \le 111$ (the extension beyond Fuller's
  computation), every optimal $n$-bit string has separator sequence
  beginning $(1, 2, 1, 1, 1, 1, 1, 1)$.  The same pattern holds for
  all best heuristic strings found through $n = 200$.
\end{observation}

Theorem~\ref{thm:sep2-collisions} gives a rigorous obstruction to
equal adjacent separators: equal consecutive separators force
$\min(b_i, b_{i+1}) \cdot \min(b_{i+1}, b_{i+2})$ collisions.
For the first pair in the computed optima, these collisions are large
enough to justify lengthening $s_2$ to~$2$; for later pairs, the
adjacent blocks are too small to warrant the observed cost.

\begin{conjecture}\label{conj:seps}
  For every $n > 7$, every optimal $n$-bit string has first separator
  length~$1$.  Moreover, for all sufficiently large~$n$, every optimal
  $n$-bit string has separator sequence beginning
  $(1, 2, 1, 1, 1, 1, 1, 1, \ldots)$.
\end{conjecture}

\subsection{Open problems}

\begin{enumerate}[label=(\roman*)]
  \item \textbf{Error term.}  The best proven bound gives
    $a(n)/n^2 = 1/2 + O(n^{-1/2})$.  Can the error term be improved
    to $O(n^{-\alpha})$ for some $\alpha > 1/2$?  Computational data
    suggest $a(n) = n^2/2 - \Theta(n^{3/2})$.

  \item \textbf{Separator structure.}  Prove or disprove
    Conjecture~\ref{conj:seps}.

  \item \textbf{Second differences.}  Prove $\DDa(n) \le 2$
    (Conjecture~\ref{conj:dda}).

  \item \textbf{Leading block.}  The empirical leading one-block
    length satisfies $b_1 \approx 0.19\,n$.  The best proven upper
    bound is $b_1 = O(n^{3/4})$
    (Corollary~\ref{cor:Lbound}).  Close this gap.

  \item \textbf{Density.}  The empirical ones-density is
    $\approx 0.85$ for moderate~$n$.  The best proven lower bound is
    $1 - O(n^{-1/4})$ (Proposition~\ref{prop:optimal-density}).
    Prove a bound approaching $1 - O(n^{-1/2})$.

  \item \textbf{Exact computation.}  Extend the certified exact range
    beyond $n = 111$.

  \item \textbf{Alphabet generalisation.}
    For an alphabet $\{0,\ldots,q{-}1\}$, define $a_q(n)$ analogously.
    We have $a_q(n)/n^2 \to 1/2$ for every $q \ge 2$: the upper
    bound $a_q(n) \le n(n{+}1)/2$ is alphabet-independent, and
    restricting to the sub-alphabet $\{0,\,q{-}1\}$ reduces to
    the binary case, giving
    $a_q(n) \ge a_2(n) \ge n^2/2 - O(n^{3/2})$.
    The deficit $n(n{+}1)/2 - a_q(n)$ appears to be much smaller for
    $q \ge 3$ (roughly $O(n \log n)$ versus $\Theta(n^{3/2})$ for
    $q = 2$), since short-substring collisions vanish once the alphabet
    exceeds~$2$.  Determining the precise deficit order for $q \ge 3$
    remains open.
\end{enumerate}

% =======================================================================
\section*{Acknowledgements}
% =======================================================================

The author thanks Martin Fuller for the original A$^*$ certification of
$a(n)$ for $n \le 80$~\cite{oeis}.  All computations were performed with
Python~3.13 using the \texttt{numpy}~\cite{numpy},
\texttt{numba}~\cite{numba}, and
\texttt{pydivsufsort}~\cite{pydivsufsort} libraries.

% =======================================================================
\begin{thebibliography}{9}

\bibitem{oeis}
  OEIS Foundation Inc.,
  \emph{The On-Line Encyclopedia of Integer Sequences},
  Sequence A112509,
  \url{https://oeis.org/A112509}, 2026.
  Values for $n \le 80$ certified by an A$^*$ search of M.~Fuller and
  updated in the OEIS log in 2025.

\bibitem{kasai2001}
  T.~Kasai, G.~Lee, H.~Arimura, S.~Arikawa, and K.~Park,
  \emph{Linear-time longest-common-prefix computation in suffix
  arrays and its applications},
  in Proceedings of CPM 2001, LNCS 2089, Springer, 2001,
  pp.~181--192.

\bibitem{pydivsufsort}
  L.~Abraham,
  \texttt{pydivsufsort}: Python bindings for DivSufSort,
  \url{https://github.com/louisabraham/pydivsufsort}, 2023.

\bibitem{numpy}
  C.~R.~Harris et al.,
  \emph{Array programming with NumPy},
  Nature \textbf{585} (2020), 357--362.

\bibitem{numba}
  S.~K.~Lam, A.~Pitrou, and S.~Seibert,
  \emph{Numba: A LLVM-based Python JIT compiler},
  Proceedings of LLVM-HPC 2015, ACM, 2015.

\bibitem{github}
  B.~Clare,
  \emph{A112509: computational exploration},
  \url{https://github.com/axelsmagichammer/A112509}, 2026.

\end{thebibliography}

\end{document}
